\documentclass[12pt]{article}
\usepackage{sectsty,setspace}
\usepackage[top=1.87cm, bottom=1.87cm, left=1.87cm, right=1.87cm]{geometry} 
\usepackage{epstopdf}
\usepackage{amsmath,latexsym,amssymb,wasysym}
\usepackage{times}
\usepackage{fancyhdr}
\usepackage[T1]{fontenc}

\setstretch{1}
%\setlength{\baselineskip}{0.167in}
\setlength\parindent{0pt}

\pagestyle{fancy}
\fancyhf{}
\fancyhead[R]{Christophe Rouleau-Desrochers}
\fancyfoot[C]{\thepage}
\renewcommand{\headrulewidth}{0pt}

\begin{document}

% ==============================================================================
\subsection*{I --- Contributions to research and development}
% ==============================================================================
% <><><><><><><><><><><><><><><><><><><><><><><><><><><><><><><><><><><><><><><>
\textbf{a. Articles published or accepted in peer-reviewed journals} \\[1pt]
% <><><><><><><><><><><><><><><><><><><><><><><><><><><><><><><><><><><><><><><>
% --- --- --- --- --- --- --- --- --- --- --- --- --- --- --- --- --- --- --- 
% Item 1 %
% --- --- --- --- --- --- --- --- --- --- --- --- --- --- --- --- --- --- --- 
Wolkovich, E.M., \textbf{Rouleau-Desrochers, C.}, García de Cortázar-Atauri, I., Walker, M.A., \& Lacombe, T. (2025) Building a more predictive model of terroir for the Anthropocene. \textit{Plants, People, Planet}. 7(6): 1611--1620. DOI: 10.1002/ppp3.70055 (BSc work). \\
I developed the meta-analysis conceptual figure and led the global mapping analysis identifying key data gaps; I assisted with the writing and edited the methods and implications sections. \\[1pt]

% <><><><><><><><><><><><><><><><><><><><><><><><><><><><><><><><><><><><><><><>
\textbf{b. Other peer-reviewed contributions} \\[1pt]
% <><><><><><><><><><><><><><><><><><><><><><><><><><><><><><><><><><><><><><><>
% --- --- --- --- --- --- --- --- --- --- --- --- --- --- --- --- --- --- --- 
% Item 1 %
% --- --- --- --- --- --- --- --- --- --- --- --- --- --- --- --- --- --- --- 
\textbf{Rouleau-Desrochers, C.}*, Van der Meersch, V., Pederson, N., \& Wolkovich, E.M. (2026) Evidence that growing season length and tree growth are decoupled in an urban arboretum. \textit{Canadian Society for Ecology and Evolution (CSEE)}. Full oral presentation at a national-level conference (MSc work). \\[1pt]

% --- --- --- --- --- --- --- --- --- --- --- --- --- --- --- --- --- --- --- 
% Item 2 %
% --- --- --- --- --- --- --- --- --- --- --- --- --- --- --- --- --- --- --- 
\textbf{Rouleau-Desrochers, C.}*, Pederson, N., \& Wolkovich, E.M. (2025) How growing season length affects tree growth in a common garden study. \textit{EcoEvo}. Oral presentation, regional-level conference (MSc work). \\[1pt]

% <><><><><><><><><><><><><><><><><><><><><><><><><><><><><><><><><><><><><><><>
\textbf{c. Non-peer-reviewed contributions} \\[1pt]
% <><><><><><><><><><><><><><><><><><><><><><><><><><><><><><><><><><><><><><><>
% --- --- --- --- --- --- --- --- --- --- --- --- --- --- --- --- --- --- --- 
% Item 1 %
% --- --- --- --- --- --- --- --- --- --- --- --- --- --- --- --- --- --- --- 
\textbf{Rouleau-Desrochers, C.}, Van der Meersch, V., Buonaiuto, D.M., Chamberlain, C.J., Loughnan, D., Pederson, N., \& Wolkovich, E.M. (2026) Warmer and longer growing seasons increase growth for juvenile trees in a common garden. \textit{Global Change Biology}. Submitted September 2026, 28 pp. (MSc work). \\
I developed the questions, collected and analysed the data, led the writing and the interpretation of the results and their implications for forest regeneration and tree responses to climate change. \\[1pt]

% --- --- --- --- --- --- --- --- --- --- --- --- --- --- --- --- --- --- --- 
% Item 2 %
% --- --- --- --- --- --- --- --- --- --- --- --- --- --- --- --- --- --- --- 
\textbf{Rouleau-Desrochers, C.}, Patry, É., Gingras, G., Ucler, J., Vézina-Doré, M., Paillé, O., \& Moghaddam, Y.Y. (2024) ``Un été parti en fumée.'' \textit{Le Point Bio}. 18: 9--19 (BSc work). \\
As lead author, I conducted the interviews with researchers and wrote the manuscript.

% ==============================================================================
\subsection*{II --- Most significant contributions to research and development}
% ==============================================================================
% --- --- --- --- --- --- --- --- --- --- --- --- --- --- --- --- --- --- --- 
% Item 1 %
% --- --- --- --- --- --- --- --- --- --- --- --- --- --- --- --- --- --- --- 
\textbf{1. Warmer and longer growing seasons increase growth for juvenile trees in a common garden (Global Change Biology and CSEE 2026).}
Extended growing seasons are widely assumed to increase tree growth, yet recent work repeatedly fails to find this link, largely because it focused on early-season leafout, when cool temperatures limit photosynthesis. I designed and led an analysis testing how juvenile tree growth in four Betulaceae species responds to both season length and thermal sum, using start- and end-of-season phenology data across four provenances. I found that warmer seasons were a better predictor of growth than season length in most species, and that end-of-season phenology was a stronger predictor than start-of-season, challenging assumptions about when in the season growth actually happens. I also tested for provenance and found little evidence of local adaptation, suggesting these responses may generalize across species ranges. I wrote the manuscript and led the interpretation; my co-authors advised on the modelling and on the tree-ring data collected at the Arnold Arboretum. We targeted \textit{Global Change Biology} for its readership of ecologists and modellers working on climate-driven shifts in ecosystem processes. My work shows that season-length effects are species- and timing-dependent, with direct implications for how carbon models incorporate phenology.

% --- --- --- --- --- --- --- --- --- --- --- --- --- --- --- --- --- --- ---
% Item 2 %
% --- --- --- --- --- --- --- --- --- --- --- --- --- --- --- --- --- --- ---
\textbf{2. Predictive model of terroir under climate change (Plants, People, Planet, 2025).}
Current vine-planting choices and sparse genotype $\times$ environment data undermine adaptation planning under climate change. I synthesized fragmented genotype $\times$ environment datasets into a proposed meta-analysis framework and designed a global data-gap map that reframed the paper's core argument: predictive models must explicitly account for missing data to guide adaptation. Working with collaborators in France and the United States, I shaped the manuscript's conceptual structure and contributed to the writing. We chose \textit{Plants, People, Planet} for its mandate to bridge fundamental plant science and applied practitioner audiences, matching our goal of making the findings usable for viticulture planning.

% ==============================================================================
\subsection*{III --- Applicant's statement}
% ==============================================================================

\textbf{Research experience}:
I first fell into natural-science research in 2021 with Dr. Pierre Drapeau, investigating woodpecker damage to electrical infrastructure across Québec, where I learned rigorous field protocols, database management and scientific coding. A 2022 Undergraduate Student Research Award (USRA) took me to remote boreal forest for fieldwork on Pileated Woodpecker habitat---banding, nest-searching and tree-classification---and I leveraged 20 years of nesting data to test how their nesting phenology is shifting with climate change, my first independent statistical project. A second USRA at UBC in 2023 (supervised by Drs. Frederik Baumgarten and E.M. Wolkovich) introduced me to large-scale experiments: over that summer I designed and implemented a defoliation experiment to test tree resilience, maintained a 1,000-tree trial, deployed dendrometers and analysed growth responses to treatments.

In 2024, I conceived and led a seven-species greenhouse experiment on growing season length, formulating the hypotheses, designing the treatments and integrating climate and phenological data into statistical models. Alongside this, I joined field campaigns in Smithers (BC) and Mount Rainier (WA) to study forest regeneration dynamics, where I improved my forestry methods and learned the basics of dendrochronology. In 2025, while continuing that experiment, I led the fieldwork at Smithers and visited the Arnold Arboretum of Harvard University to collect tree cores, deepening my expertise in dendrochronological methods.

\textbf{Relevant activities}:
Before my current academic path, I earned diplomas in hospitality and sommellerie and spent nearly ten years in those industries, where I built a strong work ethic and sense of rigour. The transition to academia began during six months working a vineyard in New Zealand, where I realized my ties to nature ran deeper than I expected. Thereafter, I enrolled in a Bachelor's degree in biology, majoring in ecology. During the first year of my MSc, I deepened my quantitative skills through a biomathematics class, a three-week workshop on hierarchical Bayesian models and the Statistical Rethinking course, refining my models in a biweekly Bayesian working group.

My work during my undergraduate studies at Université du Québec à Montréal was recognized with the Governor General's Academic Silver Medal, the Dean's List of honours and 13 competitive scholarships totalling \$76,750, alongside a perfect GPA of 4.30/4.30. For my MSc, I received 3 competitive scholarships totalling \$28,750, an honourable mention for my EcoEvo talk, and an Honours designation for my MSc thesis. I have maintained active outreach throughout: I was interviewed at the 2023 UQAM scholarship ceremony, featured by BioChambers for my growth-chamber work, hosted an open house at CSEE (2024) and presented my work twice to the Treespotters community scientists at Harvard (2025--2026).

Across the two years leading my experiment, I trained and supervised 15 undergraduate students in growth-chamber protocols, phenological monitoring and greenhouse management; 2 of them are now pursuing graduate studies. I also led the Smithers field team in 2025. I was a teaching assistant for FRST303 (2024--2025) and FRST304 (2025--2026), running office hours and preparing supplementary materials to help students grasp ecological and plant physiology concepts. Since 2021 I have also mentored 12 international students during their first months in Canada, tutored 7 peers in biochemistry and statistics, and volunteered with Biology Undergraduate Diversity in Research to advise 8 students on applying to graduate school.

My wildlife photography extends this outreach: I won the 2026 Biodiversity Research Center (UBC) photo competition, and my photos have run on the cover of \textit{Le Point Bio} (2024), in the \textit{Proceedings of the National Academy of Sciences} (2025) and in \textit{Le Rappel} (2025).

\end{document}
